\begin{mistake}{2026-07-20}{05}

\begin{problemcontent}
设 \( f(x) = (x^2 - 4x + 3)^n \)，求 \( f^{(n)}(1) \)。
\end{problemcontent}

\begin{solutioncontent}
\textbf{方法一：莱布尼茨公式}
将底数因式分解得 \( f(x) = (x-1)^n (x-3)^n \)。
设 \( u(x) = (x-1)^n \)，\( v(x) = (x-3)^n \)，应用莱布尼茨公式：
\[ f^{(n)}(x) = \sum_{k=0}^{n} C_n^k \left[ (x-1)^n \right]^{(k)} \left[ (x-3)^n \right]^{(n-k)} \]
代入 \( x=1 \)，由于 \( \left[ (x-1)^n \right]^{(k)} \big|_{x=1} = 0 \)（当 \( k < n \) 时），故展开式中仅 \( k=n \) 的项非零：
\[ f^{(n)}(1) = C_n^n \cdot n! \cdot (1-3)^n = (-2)^n n! \]

\textbf{方法二：泰勒展开系数唯一性（秒杀法）}
由 \( f(x) = (x-1)^n (x-3)^n \)，在 \( x \to 1 \) 时：
\[ \lim_{x \to 1} \frac{f(x)}{(x-1)^n} = \lim_{x \to 1} (x-3)^n = (-2)^n \]
即 \( f(x) = (-2)^n (x-1)^n + o\left((x-1)^n\right) \)。
根据带有皮亚诺余项的泰勒公式，展开式中 \( (x-1)^n \) 项的系数必然等于 \( \frac{f^{(n)}(1)}{n!} \)。
对比系数可得：
\[ \frac{f^{(n)}(1)}{n!} = (-2)^n \implies f^{(n)}(1) = (-2)^n n! \]
\end{solutioncontent}

\mistakeknowledge{
\begin{itemize}
  \item \textbf{核心公式/定理}：莱布尼茨公式；泰勒展开式的系数唯一性。
  \item \textbf{高频秒杀结论}：若 \( f(x) = (x-x_0)^n g(x) \)，且 \( g(x) \) 在 \( x_0 \) 处有直到 \( n \) 阶导数，则 \( f^{(n)}(x_0) = n! \cdot g(x_0) \)。
  \item \textbf{解题技巧}：求函数在某一点的高阶导数时，若能提取出 \( (x-x_0)^n \) 的因子，优先考虑泰勒公式系数对比法，可避免繁琐求导。
\end{itemize}}

\end{mistake}
